\chapter{Keel Mori Theorem}

The proof of the Keel-Mori theorem follows the Stacks Project. We proceed by first constructing moduli spaces on a special local class of algebraic stacks, called well-nigh affine stacks, and then gluing these local moduli spaces along
\etale\ overlaps.

\section{Local constructions}

\begin{definition}[Well-nigh]
	(\cite[Tag 0DUL][]{stacks-project})
	Let $\mathcal{X}$ be an algebraic stack. We call $\mathcal{X}$ well-nigh affine, if there exists an affine scheme $U$ and a surjective, flat, finite and finitely presented morphism $U\to \mathcal{X}$.	
\end{definition}

\begin{remark}
	The condition of being well-nigh affine is not part of the definition of an algebraic stack, but an additional finiteness and locality assumption on its presentation.
	
	Indeed, by definition of an algebraic stack (e.g.\ \cite[Tag 026N]{stacks-project}), there already exists a smooth surjective morphism $V \to \mathcal{X}$ from a scheme $V$. However, in general $V$ need neither be affine nor satisfy strong finiteness properties over $\mathcal{X}$.
	
	The well-nigh condition strengthens this by requiring the existence of a presentation $U \to \mathcal{X}$ where $U$ is affine and the morphism is flat, finite, and finitely presented. In particular, the associated equivalence relation
	\[
	R = U \times_{\mathcal{X}} U
	\]
	acquires good finiteness properties, which allows one to treat it similarly to a finite equivalence relation on an affine scheme.
	
	This additional structure is crucial in the construction of coarse moduli spaces, as it provides a local situation in which quotients can be formed and then glued via descent.
\end{remark}

\begin{lemma}[{\cite[Tag 0DUM]{stacks-project}}]
	\label{lem-groupoid-representation-of-well-nigh-affine}
	Let $\mathcal{X}$ be an algebraic stack. Then the following are equivalent:
	\begin{enumerate}
		\item $\mathcal{X}$ is well-nigh affine;
		\item there exists a groupoid scheme $(U,R,s,t,c)$ with $U$ and $R$ affine and
		$s,t \colon R \to U$ finite locally free such that
		$\mathcal{X} \cong [U/R]$.
	\end{enumerate}
	If these conditions hold, then $\mathcal{X}$ is quasi-compact, quasi-DM, and separated.
\end{lemma}


\begin{lemma}[Affine invariant quotient construction]\label{lem:affine-invariant-quotient}
	Let
	\[
	\mathcal X=[U/R]
	\]
	be a quotient stack, where 
	% https://q.uiver.app/#q=WzAsMixbMCwwLCJSPVxcU3BlY29me0J9Il0sWzEsMCwiVT1cXFNwZWNvZntBfSJdLFswLDEsInMiLDAseyJvZmZzZXQiOi0xfV0sWzAsMSwidCIsMix7Im9mZnNldCI6MX1dXQ==
	\[\begin{tikzcd}
		{R=\Specof{B}} & {U=\Specof{A}}
		\arrow["s", shift left, from=1-1, to=1-2]
		\arrow["t"', shift right, from=1-1, to=1-2]
	\end{tikzcd}\]
	is a groupoid in affine schemes. Let
	\[
	C=\{a\in A \mid s^\#(a)=t^\#(a)\}.
	\]
	Then the canonical morphism
	\[
	U\to \Specof{C}
	\]
	is $R$-invariant and therefore induces a morphism
	\[
	\mathcal X \to \Specof{C}.
	\]
\end{lemma}

\begin{proof}
	The condition
	\[
	C=\{a\in A \mid s^\#(a)=t^\#(a)\}
	\]
	means precisely that $C$ is the equalizer
	\[
	C=\operatorname{Eq}\big(A \rightrightarrows B\big),
	\]
	where the two maps are $s^\#$ and $t^\#$ corresponding to the morphisms
	\[
	s,t : R \to U.
	\]
	
	The inclusion $C \hookrightarrow A$ induces a morphism of affine schemes
	\[
	U=\Spec(A) \longrightarrow \Spec(C).
	\]
	
	We claim that this morphism is $R$-invariant. By definition, this means
	that the two compositions
	\[
	R \rightrightarrows U \to \Spec(C)
	\]
	coincide.
	
	On the level of rings, these two morphisms correspond to the two maps
	\[
	C \hookrightarrow A \xrightarrow{s^\#} B
	\qquad \text{and} \qquad
	C \hookrightarrow A \xrightarrow{t^\#} B.
	\]
	But by definition of $C$, we have $s^\#(a)=t^\#(a)$ for all $a\in C$.
	Hence these two ring maps agree, and therefore the two morphisms
	\[
	R \rightrightarrows U \to \Spec(C)
	\]
	are equal. This proves that $U \to \Spec(C)$ is $R$-invariant.
	
	It remains to show that such an invariant morphism descends to the quotient
	stack $\mathcal X=[U/R]$.
	
	By definition of the quotient stack, a morphism
	\[
	\mathcal X \to Y
	\]
	to a scheme $Y$ is equivalent to the following data:
	\begin{enumerate}
		\item a morphism $U \to Y$,
		\item such that the two pullbacks along $s,t : R \to U$ agree.
	\end{enumerate}
	
	Since we have already shown that the morphism $U \to \Spec(C)$ satisfies
	this invariance condition, it follows that it defines a morphism
	\[
	\pi : \mathcal X \to \Spec(C).
	\]
\end{proof}


\begin{lemma}[Local existence of a coarse moduli space on well-nigh affine chart]\label{lem:local-ex-cms-in-affines}
	(\cite[Tag 0DUP][]{stacks-project})
	Let $\mathcal{X}$  be a well-nigh affine algebraic stack. Then there exists a coarse moduli space $\pi:\mathcal{X} \to X$ in the category of affine schemes. Moreover $\pi$ is separated, quasi-compact and a universal homeomorphism.
\end{lemma}

\begin{proof}
	By \ref{lem-groupoid-representation-of-well-nigh-affine} there exists a presentation
	\[
	\mathcal X = [U/R],
	\]
	where
	\[
	U=\Spec(A)
	\]
	is an affine scheme and
	\[
	R = U \times_{\mathcal X} U = \Spec(B)
	\]
	is a groupoid in affine schemes such that the morphisms
	\[
	s,t : R \to U
	\]
	are finite locally free.
	
	Let
	\[
	C := \{ a \in A \mid s^\#(a) = t^\#(a) \}
	\]
	be the ring of invariant functions. Equivalently,
	\[
	C = \operatorname{Eq}\big(A \rightrightarrows B\big).
	\]
	
	By Lemma~\ref{lem:affine-invariant-quotient}, the canonical morphism
	\[
	U \to \Spec(C)
	\]
	is $R$-invariant and therefore induces a morphism $\pi: \mathcal X \to \Spec(C).$	
	Set
	\[
	X := \Spec(C).
	\]
	
	It remains to show that $\pi$ is a coarse moduli space and satisfies the
	stated properties. 
	
	By lemma \ref{lem-groupoid-representation-of-well-nigh-affine} $\mathcal{X}$ is separated. Since $X$ is affine, hence separated, $\pi$ is separated.
	$U\to \mathcal{X}$ is surjective and $U\to X$ is quasi-compact, hence $\pi$ is quasi-compact by \cite[Tag 050X]{stacks-project}.
	The details about $\pi$ being universially surjective is prooven in more detail in the proof here \cite[Tag 0DUP]{stacks-project}.
	In particular, $X=\Spec(C)$ is the desired coarse moduli space of
	$\mathcal X$ in the affine case.
\end{proof}

\begin{lemma}\label{lem:cms-pullback}
	(\cite[Tag 0DUQ][]{stacks-project})
	Let $\mathcal{X}, \mathcal{X}'$ be well-nigh affine algebraic stacks and $h: \mathcal{X}' \to \mathcal{X}$ be an \etale\ morphism, which induces isomorphisms on automorphism groups. Then the following diagram is cartesian and the induced morphism $X'\to X$ is \etale.
	% https://q.uiver.app/#q=WzAsNCxbMCwwLCJcXG1hdGhjYWx7WH0nIl0sWzEsMCwiXFxtYXRoY2Fse1h9Il0sWzAsMSwiWCciXSxbMSwxLCJYIl0sWzAsMV0sWzEsM10sWzIsM10sWzAsMl0sWzAsMywiIiwxLHsic3R5bGUiOnsibmFtZSI6ImNvcm5lciJ9fV1d
	\[\begin{tikzcd}
		{\mathcal{X}'} & {\mathcal{X}} \\
		{X'} & X
		\arrow[from=1-1, to=1-2]
		\arrow[from=1-1, to=2-1]
		\arrow["\lrcorner"{anchor=center, pos=0.125}, draw=none, from=1-1, to=2-2]
		\arrow[from=1-2, to=2-2]
		\arrow[from=2-1, to=2-2]
	\end{tikzcd}\]
\end{lemma}

\begin{lemma}\label{lem:cms-universal-along-alg-spaces}
	(\cite[Tag 0DUR][]{stacks-project})
	The morphisms of lemma \ref{lem:local-ex-cms-in-affines} is universal for morphisms into algebraic spaces.
\end{lemma}

\begin{proof}
	Let \(Y\) be an algebraic space and let
	\(h:\mathcal X\to Y\) be a morphism. We construct a unique morphism
	\(g:X\to Y\) such that \(h=g\circ \pi\).
	
	Choose an étale surjection \(p:V\to Y\), where \(V\) is a scheme. Put
	\[
	R=V\times_Y V.
	\]
	Then \(Y\) is the sheaf quotient of the étale equivalence relation
	\[
	R\rightrightarrows V.
	\]
	
	Consider the base change of $h$ along the \etale\ cover $V\to Y$:
	% https://q.uiver.app/#q=WzAsNCxbMSwxLCJZIl0sWzAsMCwiXFxtYXRoY2Fse1h9X1YiXSxbMCwxLCJWIl0sWzEsMCwiXFxtYXRoY2Fse1h9Il0sWzEsMCwiIiwwLHsic3R5bGUiOnsibmFtZSI6ImNvcm5lciJ9fV0sWzEsMiwiaF9WIiwyXSxbMiwwXSxbMywwLCJoIiwyXSxbMSwzXV0=
	\[\begin{tikzcd}
		{\mathcal{X}_V} & {\mathcal{X}} \\
		V & Y
		\arrow[from=1-1, to=1-2]
		\arrow["{h_V}"', from=1-1, to=2-1]
		\arrow["\lrcorner"{anchor=center, pos=0.125}, draw=none, from=1-1, to=2-2]
		\arrow["h"', from=1-2, to=2-2]
		\arrow[from=2-1, to=2-2]
	\end{tikzcd}\]
	
	The morphism \(\mathcal X_V\to V\) has affine target $V$.
	By Lemma \ref{lem:local-ex-cms-in-affines}, it factors uniquely as
	\[
	\mathcal X_V \xrightarrow{\pi_V} X_V
	\xrightarrow{g_V} V.
	\]
	Here \(\pi_V:\mathcal{X}_V \to X_V\) is the coarse moduli space associated to \(\mathcal X_V\) by lemma \ref{lem:local-ex-cms-in-affines}.
	Now consider
	\[
	\mathcal X_R:=\mathcal X\times_Y R,
	\]
	Equivalently, there is a canonical isomorphism
	\[
	\mathcal X_R
	\cong
	\mathcal X_V\times_{\mathcal X}\mathcal X_V,
	\]
	which can be verified by a straightforward diagram chase using the following:
	% https://q.uiver.app/#q=WzAsOCxbMCwyLCJWIl0sWzEsMiwiWSJdLFszLDEsIlYiXSxbMiwxLCJSIl0sWzEsMSwiXFxtYXRoY2Fse1h9Il0sWzAsMSwiXFxtYXRoY2Fse1h9X1YiXSxbMywwLCJcXG1hdGhjYWx7WH1fViJdLFsyLDAsIlxcbWF0aGNhbHtYfV9SIl0sWzAsMV0sWzIsMV0sWzMsMl0sWzUsNF0sWzUsMF0sWzYsMl0sWzYsNF0sWzcsNV0sWzcsNl0sWzMsMV0sWzcsNF0sWzQsMV0sWzMsMF0sWzcsM11d
	\[\begin{tikzcd}
		&& {\mathcal{X}_R} & {\mathcal{X}_V} \\
		{\mathcal{X}_V} & {\mathcal{X}} & R & V \\
		V & Y
		\arrow[from=1-3, to=1-4]
		\arrow[from=1-3, to=2-1]
		\arrow[from=1-3, to=2-2]
		\arrow[from=1-3, to=2-3]
		\arrow[from=1-4, to=2-2]
		\arrow[from=1-4, to=2-4]
		\arrow[from=2-1, to=2-2]
		\arrow[from=2-1, to=3-1]
		\arrow[from=2-2, to=3-2]
		\arrow[from=2-3, to=2-4]
		\arrow[from=2-3, to=3-1]
		\arrow[from=2-3, to=3-2]
		\arrow[from=2-4, to=3-2]
		\arrow[from=3-1, to=3-2]
	\end{tikzcd}\]
	Let \(X_R\) be the algebraic space associated to \(\mathcal X_R\) by lemma \ref{lem:local-ex-cms-in-affines}, then the morphism
	\[
	\mathcal X_R\to R
	\]
	factors uniquely as
	\[
	\mathcal X_R \xrightarrow{\pi_R} X_R
	\xrightarrow{g_R} R.
	\]
	All of this fit into the folllowing diagram:
	% https://q.uiver.app/#q=WzAsOSxbMCwyLCJSIl0sWzEsMiwiViJdLFsyLDIsIlkiXSxbMiwxLCJYX1IiXSxbMywxLCJYX1YiXSxbNCwxLCJYIl0sWzEsMCwiXFxtYXRoY2Fse1h9X1IiXSxbMiwwLCJcXG1hdGhjYWx7WH1fViJdLFszLDAsIlxcbWF0aGNhbHtYfSJdLFszLDAsImdfUiIsMl0sWzQsMSwiZ19WIiwyXSxbNSwyLCJnIiwyLHsic3R5bGUiOnsiYm9keSI6eyJuYW1lIjoiZGFzaGVkIn19fV0sWzYsMCwiaF9SIiwyLHsibGFiZWxfcG9zaXRpb24iOjQwfV0sWzYsNywiIiwwLHsib2Zmc2V0IjotMX1dLFs3LDhdLFs3LDEsImhfViIsMix7ImxhYmVsX3Bvc2l0aW9uIjo0MH1dLFs2LDMsIlxccGlfUiJdLFs3LDQsIlxccGlfViJdLFs4LDUsIlxccGkiXSxbMCwxLCIiLDAseyJvZmZzZXQiOi0xfV0sWzAsMSwiIiwwLHsib2Zmc2V0IjoxfV0sWzYsNywiIiwwLHsib2Zmc2V0IjoxfV0sWzMsNCwiIiwxLHsib2Zmc2V0IjoxfV0sWzQsNV0sWzEsMl0sWzMsNCwiIiwxLHsib2Zmc2V0IjotMX1dLFs4LDIsImgiLDIseyJsYWJlbF9wb3NpdGlvbiI6NDB9XV0=
	\begin{equation}
		\label{fig:cms-prism}
				\begin{tikzcd}
				& {\mathcal{X}_R} & {\mathcal{X}_V} & {\mathcal{X}} & \\
				&& {X_R} & {X_V} & X \\
				R & V & Y
				\arrow[shift left, from=1-2, to=1-3]
				\arrow[shift right, from=1-2, to=1-3]
				\arrow["{\pi_R}", from=1-2, to=2-3]
				\arrow["{h_R}"'{pos=0.4}, from=1-2, to=3-1]
				\arrow[from=1-3, to=1-4]
				\arrow["{\pi_V}", from=1-3, to=2-4]
				\arrow["{h_V}"'{pos=0.4}, from=1-3, to=3-2]
				\arrow["\pi", from=1-4, to=2-5]
				\arrow["h"'{pos=0.4}, from=1-4, to=3-3]
				\arrow[shift right, from=2-3, to=2-4]
				\arrow[shift left, from=2-3, to=2-4]
				\arrow["{g_R}"', from=2-3, to=3-1]
				\arrow[from=2-4, to=2-5]
				\arrow["{g_V}"', from=2-4, to=3-2]
				\arrow["g"', dashed, from=2-5, to=3-3]
				\arrow[shift left, from=3-1, to=3-2]
				\arrow[shift right, from=3-1, to=3-2]
				\arrow[from=3-2, to=3-3]
			\end{tikzcd}
	\end{equation}
	Our goal is to define $g:X\to Y$.
	The morphisms $X_R\rightrightarrows X_V$ are \etale\ and induced by
	$\mathcal{X}_R \rightrightarrows \mathcal{X}_V \to X_V$
	via lemma \ref{lem:local-ex-cms-in-affines}. By lemma \ref{lem:cms-pullback} the following squares of the diagram is cartesian:
	% https://q.uiver.app/#q=WzAsNixbMCwwLCJcXG1hdGhjYWx7WH1fUiJdLFsxLDAsIlxcbWF0aGNhbHtYfV9WIl0sWzAsMSwiWF9SIl0sWzEsMSwiWF9WIl0sWzIsMCwiXFxtYXRoY2Fse1h9Il0sWzIsMSwiWCJdLFswLDEsInRfXFxtYXRoY2Fse1h9IiwyLHsib2Zmc2V0IjoxfV0sWzEsM10sWzIsMywidF9YIiwyLHsib2Zmc2V0IjoxfV0sWzAsMl0sWzAsMywiIiwxLHsic3R5bGUiOnsibmFtZSI6ImNvcm5lciJ9fV0sWzEsNF0sWzQsNV0sWzMsNV0sWzEsNSwiIiwxLHsic3R5bGUiOnsibmFtZSI6ImNvcm5lciJ9fV0sWzAsMSwic19cXG1hdGhjYWx7WH0iLDAseyJvZmZzZXQiOi0xfV0sWzIsMywic19YIiwwLHsib2Zmc2V0IjotMX1dXQ==
	\[\begin{tikzcd}
		{\mathcal{X}_R} & {\mathcal{X}_V} & {\mathcal{X}} \\
		{X_R} & {X_V} & X
		\arrow["{t_\mathcal{X}}"', shift right, from=1-1, to=1-2]
		\arrow["{s_\mathcal{X}}", shift left, from=1-1, to=1-2]
		\arrow[from=1-1, to=2-1]
		\arrow["\lrcorner"{anchor=center, pos=0.125}, draw=none, from=1-1, to=2-2]
		\arrow[from=1-2, to=1-3]
		\arrow[from=1-2, to=2-2]
		\arrow["\lrcorner"{anchor=center, pos=0.125}, draw=none, from=1-2, to=2-3]
		\arrow[from=1-3, to=2-3]
		\arrow["{t_X}"', shift right, from=2-1, to=2-2]
		\arrow["{s_X}", shift left, from=2-1, to=2-2]
		\arrow[from=2-2, to=2-3]
	\end{tikzcd}\]
	So the back and the front of the commutative prism \ref{fig:cms-prism} are cartesian.
	Let
	\[
	s,t:R\rightrightarrows V
	\qquad\text{and}\qquad
	s_X,t_X:X_R\rightrightarrows X_V
	\]
	be the two projections. We claim that
	\[
	s\circ g_R = g_V\circ s_X,
	\qquad
	t\circ g_R = g_V\circ t_X.
	\]
	Indeed, after precomposing with \(\pi_R:\mathcal X_R\to X_R\), both
	sides of the first equality give the same morphism
	\(\mathcal X_R\to V\), namely 
	\[
		\mathcal{X}_R \xrightarrow{h_R} R \xrightarrow{s} V = \mathcal{X}_R \xrightarrow{s_\mathcal{X}} \mathcal{X}_V \xrightarrow{h_V} V.
	\]
	Since
	\(\pi_R\) is universal for morphisms to affine schemes locally on \(V\),
	uniqueness gives the equality. The second equality is analogous.
	Thus \(g_V:X_V\to V\) is compatible with the descent data
	\[
	X_R\rightrightarrows X_V
	\qquad\text{and}\qquad
	R\rightrightarrows V.
	\]
	Equivalently, we have a morphism of étale equivalence relations
	\[
	(X_R\rightrightarrows X_V)
	\longrightarrow
	(R\rightrightarrows V).
	\]
	Since
	\[
	X \cong X_V/X_R
	\qquad\text{and}\qquad
	Y \cong V/R
	\]
	as sheaves for the \etale\ topology, by the universal property of the quotient sheaf and by $Y$ being a sheaf in the \etale\ topology this morphism of equivalence relations descends to a unique morphism
	\[
	g:X\to Y.
	\]
	In fact, this is precisely indicated via Yoneda in the equalizer diagram:
	% https://q.uiver.app/#q=WzAsMyxbMCwwLCJZKFgpIl0sWzEsMCwiWShYX1YpIl0sWzIsMCwiWShYX1IpIl0sWzAsMV0sWzEsMiwiIiwwLHsib2Zmc2V0IjotMX1dLFsxLDIsIiIsMCx7Im9mZnNldCI6MX1dXQ==
	\[\begin{tikzcd}
		{Y(X)} & {Y(X_V)} & {Y(X_R)}
		\arrow[from=1-1, to=1-2]
		\arrow[shift left, from=1-2, to=1-3]
		\arrow[shift right, from=1-2, to=1-3]
	\end{tikzcd}\]
	
	It remains to check that \(h=g\circ \pi\). This can be checked after the
	étale surjection \(\mathcal X_V\to \mathcal X\). After pulling back to
	\(\mathcal X_V\), the equality becomes
	\[
	\mathcal X_V \xrightarrow{\pi_V} X_V
	\xrightarrow{g_V} V
	\longrightarrow Y
	=
	\mathcal X_V\longrightarrow V\longrightarrow Y,
	\]
	which holds by construction of \(g_V\). Hence \(h=g\circ \pi\).
	
	Finally, suppose \(g,g':X\to Y\) both satisfy
	\[
	g\circ \pi = h = g'\circ \pi.
	\]
	Pulling back along the étale cover \(X_V\to X\), the two morphisms
	\(g|_{X_V}\) and \(g'|_{X_V}\) both factor the same morphism
	\(\mathcal X_V\to V\). By the uniqueness part of
	Lemma \ref{lem:local-ex-cms-in-affines}, they agree on \(X_V\).
	Since \(X_V\to X\) is an étale surjection, it follows that
	\(g=g'\).
	
	Thus \(f:\mathcal X\to X\) is universal for morphisms from
	\(\mathcal X\) to any algebraic space.
\end{proof}


\begin{lemma}\label{lem:cms-pullback-domain-non-well-nigh-affine}
	(\cite[Tag 0DUS][]{stacks-project})
	Let $h:\mathcal{X'} \to \mathcal{X}$ be an \etale\ morphism of algebraic stacks, which induces isomorphism on automorphism groups and assume further that $\mathcal{X}$ is well-nigh affine. Then there exists a cartesian diagram 
	% https://q.uiver.app/#q=WzAsNCxbMCwwLCJcXG1hdGhjYWx7WH0nIl0sWzEsMCwiXFxtYXRoY2Fse1h9Il0sWzAsMSwiWCciXSxbMSwxLCJYIl0sWzAsMl0sWzAsMSwiaCIsMl0sWzEsM10sWzIsM10sWzAsMywiIiwxLHsic3R5bGUiOnsibmFtZSI6ImNvcm5lciJ9fV1d
	\[\begin{tikzcd}
		{\mathcal{X}'} & {\mathcal{X}} \\
		{X'} & X,
		\arrow["h"', from=1-1, to=1-2]
		\arrow[from=1-1, to=2-1]
		\arrow["\lrcorner"{anchor=center, pos=0.125}, draw=none, from=1-1, to=2-2]
		\arrow[from=1-2, to=2-2]
		\arrow[from=2-1, to=2-2]
	\end{tikzcd}\]
	such that $M'\to M$ is \etale\ and $\mathcal{X}\to X$ is a coarse moduli space.
\end{lemma}

So far we have prooven the existence of a coarse moduli space for well-nigh affine stacks. In the next part, we will construct for an algebraic stack with finite inertia a cover of such well-nigh affine stacks.


\begin{lemma}\label{lem:wna-decomposition}
	(\cite[Tag 0DUE]{stacks-project})
	Let $\mathcal{X}$ be an algebraic stack with finite inertia. Then $\mathcal{X}$ admits an \etale\ cover by well-nigh affine stacks.
	
	More precisely, there exists a family of morphisms
	\[
	g_i : \mathcal{X}_i \to \mathcal{X}
	\]
	such that:
	\begin{enumerate}
		\item the morphisms $g_i$ are representable by algebraic spaces, separated, and étale,
		
		\item the induced map
		\[
		\coprod_i |\mathcal{X}_i| \to |\mathcal{X}|
		\]
		is surjective,
		
		\item each $\mathcal{X}_i$ is well-nigh affine,
		
		\item the natural morphism
		\[
		\mathcal{I}_{\mathcal{X}_i}
		\longrightarrow
		\mathcal{X}_i \times_{\mathcal{X}} \mathcal{I}_{\mathcal{X}}
		\]
		is an isomorphism.
	\end{enumerate}
\end{lemma}


\begin{theorem}[Keel-Mori Theorem]
	Let $\mathcal{X}$ be an algebraic stack with finite inertia, i.e. $\mathcal{I_X}\to \mathcal{X}$ is a finite morphism. Then there exists a coarse moduli space
	\[
	\pi: \mathcal{X}\to X.
	\]
\end{theorem}

\begin{proof}
	By lemma \ref{lem:wna-decomposition} there is a decomposition into well-nigh affine substacks $g_i:\mathcal{X}_i\to \mathcal{X}$. By  \ref{lem:local-ex-cms-in-affines} and \ref{lem:cms-universal-along-alg-spaces} there exists coarse moduli spaces $\pi_i:\mathcal{X}_i \to X_i$. We now have to ensure, that these constructed local coarse moduli spaces glue to a global coarse moduli space $X$.
	Therefore let $\mathcal{X}_{ij}$ be the pullback
	% https://q.uiver.app/#q=WzAsNCxbMCwwLCJcXG1hdGhjYWx7WH1fe2lqfSJdLFsxLDAsIlxcbWF0aGNhbHtYfV9pIl0sWzAsMSwiXFxtYXRoY2Fse1h9X2oiXSxbMSwxLCJcXG1hdGhjYWx7WH0iXSxbMiwzLCJnX2oiLDJdLFsxLDMsImdfaSJdLFswLDFdLFswLDJdLFswLDMsIiIsMSx7InN0eWxlIjp7Im5hbWUiOiJjb3JuZXIifX1dXQ==
	\[\begin{tikzcd}
		{\mathcal{X}_{ij}} & {\mathcal{X}_i} \\
		{\mathcal{X}_j} & {\mathcal{X}}
		\arrow[from=1-1, to=1-2]
		\arrow[from=1-1, to=2-1]
		\arrow["\lrcorner"{anchor=center, pos=0.125}, draw=none, from=1-1, to=2-2]
		\arrow["{g_i}", from=1-2, to=2-2]
		\arrow["{g_j}"', from=2-1, to=2-2]
	\end{tikzcd}\]
	By lemma \ref{lem:wna-decomposition} the morphisms $g_i$ are representable by algebraic spaces and \etale. Since \etale\ morphisms are stable under base change (\cite[see][Tag 0CIN]{stacks-project}) , all $\mathcal{X}_{ij}\to \mathcal{X}_i$ are \etale. All $\mathcal{X}_{ij}\to \mathcal{X}_i$ induce isomorphisms on automorphisms groups, seen by the following diagram:
	% https://q.uiver.app/#q=WzAsOCxbMCwyLCJcXG1hdGhjYWx7WH1faiJdLFsyLDEsIlxcbWF0aGNhbHtJfV9cXG1hdGhjYWx7WF9pfSJdLFsxLDIsIlxcbWF0aGNhbHtYfSJdLFszLDEsIlxcbWF0aGNhbHtJX1h9Il0sWzMsMCwiXFxtYXRoY2Fse0l9X1xcbWF0aGNhbHtYX2l9Il0sWzIsMCwiXFxtYXRoY2Fse0l9X3tcXG1hdGhjYWx7WH1fe2lqfX0iXSxbMSwxLCJcXG1hdGhjYWx7WH1faSJdLFswLDEsIlxcbWF0aGNhbHtYfV97aWp9Il0sWzMsMl0sWzEsMF0sWzEsM10sWzAsMl0sWzEsMiwiIiwxLHsic3R5bGUiOnsibmFtZSI6ImNvcm5lciJ9fV0sWzUsNF0sWzQsM10sWzUsMV0sWzcsMF0sWzcsNl0sWzYsMl0sWzcsMiwiIiwxLHsic3R5bGUiOnsibmFtZSI6ImNvcm5lciJ9fV0sWzQsNl0sWzUsN10sWzQsMiwiIiwxLHsic3R5bGUiOnsibmFtZSI6ImNvcm5lciJ9fV1d
	\[\begin{tikzcd}
		&& {\mathcal{I}_{\mathcal{X}_{ij}}} & {\mathcal{I}_\mathcal{X_i}} \\
		{\mathcal{X}_{ij}} & {\mathcal{X}_i} & {\mathcal{I}_\mathcal{X_i}} & {\mathcal{I_X}} \\
		{\mathcal{X}_j} & {\mathcal{X}}
		\arrow[from=1-3, to=1-4]
		\arrow[from=1-3, to=2-1]
		\arrow[from=1-3, to=2-3]
		\arrow[from=1-4, to=2-2]
		\arrow[from=1-4, to=2-4]
		\arrow["\lrcorner"{anchor=center, pos=0.125, rotate=-90}, draw=none, from=1-4, to=3-2]
		\arrow[from=2-1, to=2-2]
		\arrow[from=2-1, to=3-1]
		\arrow["\lrcorner"{anchor=center, pos=0.125}, draw=none, from=2-1, to=3-2]
		\arrow[from=2-2, to=3-2]
		\arrow[from=2-3, to=2-4]
		\arrow[from=2-3, to=3-1]
		\arrow["\lrcorner"{anchor=center, pos=0.125, rotate=-90}, draw=none, from=2-3, to=3-2]
		\arrow[from=2-4, to=3-2]
		\arrow[from=3-1, to=3-2]
	\end{tikzcd}\]
	Since the front, the bottom and right squares are cartesian, the left back face is cartesian.
	In general, the pullback of two well-nigh affine stacks is not again well-nigh affine. Hence, we can only apply lemma \ref{lem:cms-pullback-domain-non-well-nigh-affine} here.
	Now since $\mathcal{X}_i$ is well-nigh affine, by lemma \ref{lem:cms-pullback-domain-non-well-nigh-affine} applied to $\mathcal{X}_{ij} \to \mathcal{X}_i$ we get cartesian squares
	\begin{equation}
		\label{diag:coarse-moduli-cartesian-square}
		% https://q.uiver.app/#q=WzAsNixbMSwwLCJcXG1hdGhjYWx7WH1fe2lqfSJdLFsyLDAsIlxcbWF0aGNhbHtYfV9pIl0sWzAsMCwiXFxtYXRoY2Fse1h9X2oiXSxbMCwxLCJYX2oiXSxbMSwxLCJYX3tpan0iXSxbMiwxLCJYX2kiXSxbMCwxXSxbMCwyXSxbMSw1LCJcXHBpX2kiXSxbMCw0LCJcXHBpX3tpan0iXSxbMiwzLCJcXHBpX2oiXSxbNCwzXSxbNCw1XSxbMCw1LCIiLDAseyJzdHlsZSI6eyJuYW1lIjoiY29ybmVyIn19XSxbMCwzLCIiLDAseyJzdHlsZSI6eyJuYW1lIjoiY29ybmVyIn19XV0=
		\begin{tikzcd}
			{\mathcal{X}_j} & {\mathcal{X}_{ij}} & {\mathcal{X}_i} \\
			{X_j} & {X_{ij}} & {X_i}
			\arrow["{\pi_j}", from=1-1, to=2-1]
			\arrow[from=1-2, to=1-1]
			\arrow[from=1-2, to=1-3]
			\arrow["\lrcorner"{anchor=center, pos=0.125, rotate=-90}, draw=none, from=1-2, to=2-1]
			\arrow["{\pi_{ij}}", from=1-2, to=2-2]
			\arrow["\lrcorner"{anchor=center, pos=0.125}, draw=none, from=1-2, to=2-3]
			\arrow["{\pi_i}", from=1-3, to=2-3]
			\arrow[from=2-2, to=2-1]
			\arrow[from=2-2, to=2-3]
		\end{tikzcd}
	\end{equation}

	with $X_{ij}\to X_i$ \etale\ and $\pi_i:\mathcal{X}_i\to X_i$ coarse moduli spaces. We now glue the $X_i$ to a global coarse moduli space $X$ via \etale\ descent. Claim:
	% https://q.uiver.app/#q=WzAsMixbMCwwLCJcXGNvcHJvZCBYX3tpan0iXSxbMSwwLCJcXGJpZyhcXGNvcHJvZCBYX2lcXGJpZylcXHRpbWVzIFxcYmlnKFxcY29wcm9kIFhfalxcYmlnKSJdLFswLDEsIihzLHQpIl1d
	\[\begin{tikzcd}
		{\coprod X_{ij}} & {\big(\coprod X_i\big)\times \big(\coprod X_j\big)}
		\arrow["{(s,t)}", from=1-1, to=1-2]
	\end{tikzcd}\]
	is an \etale\ equivalence relation, where $s$ and $t$ are the induces morphisms from \ref{diag:coarse-moduli-cartesian-square}.
	
	Proof: We know that $s, t$ are \etale, since all $X_{ij} \to X_i$ are \etale. Setting $R=\coprod X_{ij}$ and $U=\coprod X_i$, diagram \ref{diag:coarse-moduli-cartesian-square} translates to 
	\[
	R \to U\times U.
	\]
	We now construct a groupoid structure on $\coprod X_i$ as in definition \ref{def:groupoid-in-alg-spaces}, in order to define the quotient algebraic space $U/R$.
	
	By associativity of the fibre product, we have 
	\[
	\mathcal{X}_{ijk}:= \mathcal{X}_i\times_\mathcal{X} \mathcal{X}_j \times_\mathcal{X} \mathcal{X}_k \cong \mathcal{X}_{ij} \times_{\mathcal{X}_j} \mathcal{X}_{jk}.
	\]
	The induced morphism $\mathcal{X}_{ijk} \to \mathcal{X}_j$ is again \etale\ and induces isomorphisms on automorphism groups as a base change along $g_i$ by lemma \ref{lem:wna-decomposition}.
	Now apply lemma \ref{lem:cms-pullback-domain-non-well-nigh-affine} again to $\mathcal{X}_{ijk} \to \mathcal{X}_i$ to obtain a cartesian square:
	% https://q.uiver.app/#q=WzAsNCxbMCwwLCJcXG1hdGhjYWx7WH1fe2lqa30iXSxbMSwwLCJcXG1hdGhjYWx7WH1faSJdLFswLDEsIlhfe2lqa30iXSxbMSwxLCJYX2kiXSxbMCwxXSxbMCwyXSxbMiwzXSxbMSwzXSxbMCwzLCIiLDEseyJzdHlsZSI6eyJuYW1lIjoiY29ybmVyIn19XV0=
	\[\begin{tikzcd}
		{\mathcal{X}_{ijk}} & {\mathcal{X}_i} \\
		{X_{ijk}} & {X_i}
		\arrow[from=1-1, to=1-2]
		\arrow[from=1-1, to=2-1]
		\arrow["\lrcorner"{anchor=center, pos=0.125}, draw=none, from=1-1, to=2-2]
		\arrow[from=1-2, to=2-2]
		\arrow[from=2-1, to=2-2]
	\end{tikzcd}\]
	This diagram factors via the following:
	% https://q.uiver.app/#q=WzAsNixbMCwwLCJcXG1hdGhjYWx7WH1fe2lqa30iXSxbMSwwLCJcXG1hdGhjYWx7WH1fe2lqfSJdLFswLDEsIlhfe2lqa30iXSxbMSwxLCJYX3tpan0iXSxbMiwwLCJcXG1hdGhjYWx7WH1faiJdLFsyLDEsIlhfaiJdLFswLDFdLFswLDJdLFsxLDNdLFsyLDNdLFsxLDRdLFszLDVdLFs0LDVdXQ==
	\[\begin{tikzcd}
		{\mathcal{X}_{ijk}} & {\mathcal{X}_{ij}} & {\mathcal{X}_j} \\
		{X_{ijk}} & {X_{ij}} & {X_j}
		\arrow[from=1-1, to=1-2]
		\arrow[from=1-1, to=2-1]
		\arrow[from=1-2, to=1-3]
		\arrow[from=1-2, to=2-2]
		\arrow[from=1-3, to=2-3]
		\arrow[from=2-1, to=2-2]
		\arrow[from=2-2, to=2-3]
	\end{tikzcd}\]
	Since the right square and the outer rectagle is cartesian, the left square is also cartesian. 
	Now consider the following commutative diagram:
	% https://q.uiver.app/#q=WzAsOCxbMSwxLCJcXG1hdGhjYWx7WH1faiJdLFswLDEsIlxcbWF0aGNhbHtYfV97amt9Il0sWzMsMCwiXFxtYXRoY2Fse1h9X3tpan0iXSxbMiwwLCJcXG1hdGhjYWx7WH1fe2lqa30iXSxbMSwyLCJYX2oiXSxbMCwyLCJYX3tqa30iXSxbMiwxLCJYX3tpamt9Il0sWzMsMSwiWF97aWp9Il0sWzEsMF0sWzMsMV0sWzIsMF0sWzMsMl0sWzMsNl0sWzIsN10sWzYsNV0sWzcsNF0sWzUsNF0sWzEsNV0sWzAsNF0sWzYsN10sWzMsMF0sWzYsNF1d
	\[\begin{tikzcd}
		&& {\mathcal{X}_{ijk}} & {\mathcal{X}_{ij}} \\
		{\mathcal{X}_{jk}} & {\mathcal{X}_j} & {X_{ijk}} & {X_{ij}} \\
		{X_{jk}} & {X_j}
		\arrow[from=1-3, to=1-4]
		\arrow[from=1-3, to=2-1]
		\arrow[from=1-3, to=2-2]
		\arrow[from=1-3, to=2-3]
		\arrow[from=1-4, to=2-2]
		\arrow[from=1-4, to=2-4]
		\arrow[from=2-1, to=2-2]
		\arrow[from=2-1, to=3-1]
		\arrow[from=2-2, to=3-2]
		\arrow[from=2-3, to=2-4]
		\arrow[from=2-3, to=3-1]
		\arrow[from=2-3, to=3-2]
		\arrow[from=2-4, to=3-2]
		\arrow[from=3-1, to=3-2]
	\end{tikzcd}\]
\end{proof}